\cotefile{103}
\scanpage{11}
\deleted{\unsure{catégories|catégorie}}\inserted{\unsure{champs}} de Picard \unsure{épinglés} par $M,N$, $\unsure{\sigma|o}$ étant l'invariant qui \inserted{s'obtient en} associant à toute section \inserted{$L$} d'un champ de Picard la symétrie de $\deleted{\unsure{K|L}}\,L\otimes L$, interprétée comme section de ${}_{2}N$. Je sais prouver (sauf erreur) que tout hom. $M\to{}_{2}N$ provient d'un champ de Picard convenable (épinglé par $M,N$) (\unsure{a|à} priori l'obstruction est dans $\operatorname{Ext}^{3}(X;M,N)$, mais un argument ``universel'' prouve qu'elle est nulle). Cela prouve que l'extension $\alpha$, est bien proche d'être splittée :\deleted{\unsure{u|un}} tte section du troisième faisceau, sur un objet quelconque de $X$, se remonte - en d'autres termes l'extension a une section ``ensembliste''. Bien sûr, il y a mieux en fait : toute section sur un $U\in\operatorname{Ob}X$ ``provient'' d'un élément de ${\unsure{\mathcal H|\mathscr H}}^{2}(U,\unsure{E'|E^{1}}(M,N))$\unsure{,|.} (hypercohomologie $\mathcal H^{2}$).
% archive sheet number: 10
